\documentclass[10pt,a4paper]{amsart}
\usepackage[margin=19mm]{geometry}
\usepackage{amsmath,amssymb,mathtools}
\usepackage[scr=rsfs,cal=euler]{mathalfa}
\usepackage{xfrac}
\usepackage[unicode,hidelinks,bookmarks=false]{hyperref}
\newcommand{\ie}{i.e.\ }
\newcommand{\cf}{cf.\ }
\newcommand{\resp}{respectively\ }
\newcommand{\Subsection}{\subsection}
\providecommand{\nobreakdash}{\nobreak\hbox{-}}
\setlength{\emergencystretch}{2em}
\multlinegap0pt
\usepackage{amssymb}
\usepackage{bm}
\usepackage{float}
\usepackage{mathtools}
\newtheorem{prop}{Proposition}
\newcommand{\R}{\mathbb{R}}
\title{On the Digits of Partition Functions}
\author{Siddharth Iyer}
\date{Source: arXiv:2602.07726v2 (CC BY 4.0)}
\begin{document}
\maketitle

\noindent\emph{Source and licence.} On the Digits of Partition Functions. Version arXiv:2602.07726v2 (CC BY 4.0), distributed by arXiv under the Creative Commons Attribution 4.0 International licence, http://creativecommons.org/licenses/by/4.0/, as recorded in the arXiv licence metadata for this exact version and shown on the abstract page https://arxiv.org/abs/2602.07726v2. Full licence notice: \texttt{licenses/arxiv-2602.07726-CC-BY-4.0.txt}.\par\smallskip

\noindent\textit{Editorial scope.} Counted unit: the article's prop prop:framework (source0.tex lines 142--157). The article's complete original proof of it is delivered with it (source0.tex lines 158--180, one \texttt{proof} environment of 23 lines). No mathematics is written, repaired or re-derived: the statement and the proof are the article's own lines, in the article's own notation, and no retained line is substituted or re-flowed.  No further source-local context is needed: every cross-reference of every retained line resolves inside the two delivered spans.  Nothing was omitted from any retained span, so the package carries no omission record.  No retained line contains a citation, so the wrapper carries no bibliography and no bibliography entry.  No retained line contains a figure environment, an image-inclusion command or drawing code, so no image is delivered and the package declares no asset dependency.  Editorial formatting only, in addition: an AMS \texttt{amsart} preamble that restates the article's own declarations and macros used by the retained lines, namely the environments \texttt{prop} on separate counters, together with the standard packages those lines need.  The retained environments print in this document as Proposition 1 in order of appearance, whereas the article prints the same items with the numbering of its own full document.  The complete native source of the article as distributed in the arXiv e-print for this exact version is bundled unchanged as \texttt{sources/194204/source0.tex}.
\begin{prop}
\label{prop:framework}
Let $g: [K, \infty) \rightarrow \R$ be a function of the form 
\begin{equation*}
g(x) = h(x) + E(x) = c_{1}x^{\theta}+ c_{2}\ln x+c_{3}+E(x),
\end{equation*}
where $c_{1} > 0$, $c_{2} \leq 0$, $c_{3}\in \R$, $0<\theta<1$, and $|E(x)| \leq c_{4}x^{-\theta}$ for some constant $c_{4}>0$. Let $\delta > 0$ and $[a, a+\delta) \subseteq [0,1)$. Define the following constants:
\begin{align*}
L_1 &= \left(\frac{-3c_{2}}{c_{1}\theta}\right)^{1/\theta}, & L_2 &= \left(\frac{3 c_{4}}{\delta}\right)^{1/\theta}, & L_3 &= \left(\frac{2}{c_1 \theta 2^{\theta-1}}\right)^{1/\theta}, \\
L_4 &= \left(\frac{3 c_{1} \theta}{\delta} \right)^{\frac{1}{1-\theta}}, & D &= \frac{2}{c_{1}\theta 2^{\theta-1}}. & &
\end{align*}
Then, there exists an integer $m$ satisfying $\{g(m)\} \in [a,a+\delta)$ such that
\begin{equation*}
m \leq 2\max\{K, L_{1}, L_{2}+1, L_{3}, L_{4}\}.
\end{equation*}
\end{prop}

\begin{proof}
By differentiating $h(x)$, we obtain $h'(x) = c_{1}\theta x^{\theta-1} + c_{2}/x$. For $x \geq L_{1}$, it is straightforward to check that $h'(x) \geq \frac{2}{3}c_{1}\theta x^{\theta-1}$. Furthermore, for $x \geq 1$, we clearly have $h'(x) \leq c_{1}\theta x^{\theta-1}$ since $c_2 \leq 0$. 

Notice that when $x \geq L_{3} = D^{1/\theta}$, we have $Dx^{1-\theta} \leq x$, which implies the interval $[x, x+Dx^{1-\theta}]$ is contained within $[x, 2x]$. For $x \geq \max\{L_1, L_3\}$ we have
\begin{align*}
h(x+Dx^{1-\theta}) - h(x) &= \int_{x}^{x+Dx^{1-\theta}} h'(y) \, dy \\
&\geq \frac{2}{3} \int_x^{x+Dx^{1-\theta}} c_1 \theta y^{\theta-1} \, dy \\
&= \frac{2}{3} c_1 \left( (x+Dx^{1-\theta})^\theta - x^\theta \right).
\end{align*}
Applying the Mean Value Theorem, the difference $(x+Dx^{1-\theta})^\theta - x^\theta$ evaluates to $\theta C^{\theta-1} Dx^{1-\theta}$ for some $C \in [x, 2x]$. Since $\theta - 1 < 0$, the minimum is attained at $C = 2x$. Thus:
\begin{equation*}
h(x+Dx^{1-\theta}) - h(x) \geq \frac{2}{3} c_1 \theta (2x)^{\theta-1} D x^{1-\theta} = \frac{2}{3} c_1 \theta 2^{\theta-1} D = \frac{4}{3}.
\end{equation*}
Because $4/3 > 1 + \delta/3$, the interval $[h(x), h(x+Dx^{1-\theta})]$ is sufficiently large to choose values $y_{1} < y_{2}$ such that $y_1 \equiv a + \delta/3 \pmod 1$ and $y_{2} = y_{1} + \delta/3$. By continuity, there exist $x_1 < x_2$ in $[x, x+Dx^{1-\theta}]$ such that $h(x_i) = y_i$. 

Using the Mean Value Theorem once more, $y_2 - y_1 = h'(c)(x_2 - x_1)$ for some $c \in [x_1, x_2]$. Because $h'(c) \leq c_1 \theta c^{\theta-1} \leq c_1 \theta x^{\theta-1}$, we find:
\begin{equation*}
x_{2}-x_{1} \geq \frac{\delta/3}{c_1 \theta x^{\theta-1}} = \frac{\delta}{3 c_{1}\theta}x^{1-\theta}.
\end{equation*}
If we require $x \geq L_{4}$, then $x_{2}-x_{1} \geq 1$. Consequently, the interval $[x_{1},x_{2}]$ must contain an integer $m$. Because $h(m) \in [y_1, y_2]$, we know $\{h(m)\} \in [a+\delta/3, a+2\delta/3]$. Finally, if $x > L_2$, we guarantee that the error term satisfies $|E(m)| \leq c_4 x^{-\theta} < \delta/3$. Therefore, $\{g(m)\} = \{h(m) + E(m)\} \in (a, a+\delta)$.

Taking $x = \max\{K, L_1, L_2+1, L_3, L_4\}$, the integer $m$ constructed above satisfies $m \in [x,x+Dx^{1-\theta}]$, and thus $m \leq 2x$.
\end{proof}

\end{document}
